\section{Formal Burnside lift and the squarefree certificate}
\label{sec:burnside}

\subsection{Squarefree cycles are cyclic set partitions}

Fix $|P|=n$ and consider the full squarefree multidegree
$x_P=\prod_{p\in P}x_p$.  Every label appears exactly once.  The edge
subsets of a contributing word are therefore disjoint and cover $P$.

\begin{proposition}[Cyclic-partition enumeration]
\label{prop:cyclic-partition}
A cyclic word at squarefree content $x_P$ is primitive and is equivalent to a
cyclically ordered set partition of $P$.  With $m$ blocks, the number of
cycles is
\[
  (m-1)!S(n,m),
  \tag{4.1}
\]
and their common scalar sign is $(-1)^{n+m}$.
\end{proposition}

\begin{proof}
A nontrivial repetition would repeat every label of the shorter root, which
is impossible at squarefree content.  An unordered set partition into $m$
blocks has $(m-1)!$ cyclic orders.  The product of the edge signs is
\[
 \prod_{j=1}^m(-1)^{|S_j|+1}=(-1)^{n+m}.
\]
\end{proof}

Consequently the total number of squarefree primitive cycles is
\[
  N_n=\sum_{m=1}^n(m-1)!S(n,m).
  \tag{4.2}
\]
For $n\ge2$, the positive and negative scalar counts agree.  One proof is to
observe that their difference is the coefficient of $x_P$ in
\[
  -\log\prod_{p\in P}(1-x_p)
  =\sum_{p\in P}-\log(1-x_p),
  \tag{4.3}
\]
which has no mixed monomial.  Exact values through seven labels are listed in
\Cref{tab:squarefree-counts}.  These counts show that scalar balance persists;
they do not prove equivariant balance.

\begin{table}[t]
\centering
\caption{Squarefree primitive cyclic-partition counts.  The equality of
positive and negative columns is scalar cancellation; it does not imply an
equivariant pairing.}
\label{tab:squarefree-counts}
\begin{tabular}{rrrrr}
\toprule
$n$ & total & positive & negative & difference\\
\midrule
2 & 2 & 1 & 1 & 0\\
3 & 6 & 3 & 3 & 0\\
4 & 26 & 13 & 13 & 0\\
5 & 150 & 75 & 75 & 0\\
6 & 1,082 & 541 & 541 & 0\\
7 & 9,366 & 4,683 & 4,683 & 0\\
\bottomrule
\end{tabular}
\end{table}

\subsection{The first nonzero Burnside residual}

Take $P=\{p,q,r\}$.  The six cycles separate by sign as
\[
 C_+=\{[pqr],[p][q][r],[p][r][q]\},
 \tag{4.4}
\]
\[
 C_-=\{[p][qr],[q][pr],[r][pq]\}.
 \tag{4.5}
\]
Here $[pqr]$ denotes the one-edge subset $\{p,q,r\}$, while brackets between
singletons denote a cyclic word of several edges.

\begin{theorem}[$pqr$ Burnside and character certificate]
\label{thm:pqr}
The virtual $S_3$-set at squarefree content $pqr$ is
\[
  \RP_3=[C_+]-[C_-]
  =[S_3/S_3]+[S_3/C_3]-[S_3/C_2].
  \tag{4.6}
\]
Its marks at subgroup classes $(1,C_2,C_3,S_3)$ are
\[
  (0,0,3,1).
  \tag{4.7}
\]
Under permutation linearization,
\[
  R_3=\one\oplus\sgnrep-\stdrep,
  \tag{4.8}
\]
and its character at $(e,(12),(123))$ is
\[
  (0,0,3).
  \tag{4.9}
\]
Thus scalar dimension kills a nonzero Burnside and representation class.
\end{theorem}

\begin{proof}
The one-block cycle in (4.4) is fixed by all of $S_3$, giving $S_3/S_3$.
The two orientations of three singleton blocks form the two-point orbit
$S_3/C_3$: a transposition exchanges them, while a three-cycle is a cyclic
rotation and fixes both necklaces.  The three singleton--pair cycles in
(4.5) form the natural orbit $S_3/C_2$.  This proves (4.6).

The identity subgroup fixes all objects.  A subgroup $C_2$ fixes the
one-block positive cycle and one singleton--pair negative cycle.  The $C_3$
subgroup fixes all three positive cycles and none of the negative cycles.
The full group fixes only the one-block positive cycle.  Subtraction gives
(4.7).

Finally,
\[
 \C[S_3/S_3]=\one,
 \quad \C[S_3/C_3]=\one\oplus\sgnrep,
 \quad \C[S_3/C_2]=\one\oplus\stdrep.
 \tag{4.10}
\]
Their virtual difference is (4.8).  The $S_3$ character table then gives
(4.9).
\end{proof}

\begin{table}[t]
\centering
\caption{The exact $pqr$ certificate.  Marks count subgroup-fixed objects;
the final column is positive minus negative.}
\label{tab:pqr-marks}
\begin{tabular}{lrrr}
\toprule
Subgroup & $|C_+^H|$ & $|C_-^H|$ & $\phi_H(\RP_3)$\\
\midrule
$1$ & 3 & 3 & 0\\
$C_2$ & 1 & 1 & 0\\
$C_3$ & 3 & 0 & 3\\
$S_3$ & 1 & 0 & 1\\
\bottomrule
\end{tabular}
\end{table}

An equivariant sign-reversing bijection would induce equal fixed-point counts
for every subgroup.  The $C_3$ mark rules it out.  The failure is stronger
than unequal presentation choices: it is invariant under relabeling and is
visible in both Burnside and representation rings.

\subsection{Adams isolation and formal projectivity}

\begin{proposition}[Squarefree Adams isolation]
\label{prop:adams-isolation}
No Adams power $\psi^k$ with $k>1$ maps a monomial of integral nonnegative
multidegree to $x_px_qx_r$.  Hence no temporal higher-power counterterm
removes $\RP_3$ or $R_3$.
\end{proposition}

\begin{proof}
On monomials, $\psi^k(x^\alpha)=x^{k\alpha}$.  Equality
$k\alpha=(1,1,1)$ has no solution in $\N^3$ when $k>1$.
\end{proof}

The $C_2$ carrier from (3.11) handles the sign component simultaneously:
$\psi^k(\tau)=\tau^k$.  Thus the formal positive result respects both
multidegree and scalar temporal powers.

For an inclusion $P\subset Q$, specialize $x_q=0$ for every
$q\in Q\setminus P$.  Every edge meeting a new label vanishes, while the old
edge and primitive-cycle terms remain.  These zero-specializations compose,
so the multigraded Burnside/species ledgers form a formal projective family.
This statement concerns coefficients and cycle indices.  It does not yet
construct an inductive system of fixed weighted operators; that analytic
question is deferred to \Cref{sec:infinite}.
